% CH6 reading draft, 8 September 2026.
% Structure follows Chapters/CH6_object_definition.tex.
% Analysis blueprint: ANA_EXOT_2025_06_INT1, 20 April 2026, Section 5.
% Paper cross-check: arXiv:2606.02067v1, Sections 5 and 7.
% Source decisions and unresolved production details: CH6_source_notes.md.
\chapter{Object Reconstruction and Identification}
\label{chapter:object_reconstruction}

Hadronically decaying $\tau$-leptons, electrons, muons and jets are used in this analysis.
This chapter describes their reconstruction, identification and calibration, followed by overlap removal and the reconstruction of $\met$.

These objects are reconstructed and identified using detector measurements.
Tracks are used to determine charged-particle momenta and production vertices, while energy deposits are measured in the calorimeters.
The reconstructed energies and momenta are calibrated to correct biases in the detector response.
Simulated event weights are further corrected using efficiency scale factors.


\section{Tracks}
\label{sec:ch6_tracks}

A track is a fitted trajectory of a charged particle through the inner detector.
The solenoidal magnetic field bends the trajectory in the transverse plane, so that its curvature determines the charge-to-momentum ratio.
Measurements close to the interaction point also determine whether the particle is compatible with production in the primary collision or originates from a displaced decay.
These two functions make tracking central to both lepton reconstruction and flavour tagging.

Track reconstruction begins by grouping neighbouring signals in the pixel and silicon microstrip detectors into clusters~\cite{ATLAS2017DenseTracking}.
Pixel clusters provide three-dimensional position measurements, called space points.
For the microstrip detector, measurements on the two sides of a module are combined to obtain a space point.
Compatible combinations of space points form initial track seeds, which are extended through additional detector layers.
During this extension, a combinatorial Kalman filter updates the trajectory estimate as measurements are added, accounting for measurement uncertainties and the effects of material along the trajectory.
Several candidates may initially use the same measurements.
An ambiguity-resolution procedure selects compatible candidates using their fit quality, hit content and missing expected measurements.
Accepted silicon tracks are extended into the transition radiation tracker where suitable measurements are available.

The fitted trajectory can be specified by five parameters at its point of closest approach to a reference beam line:
\begin{align}
    \left(d_0,\ z_0,\ \phi_0,\ \theta,\ q/p\right).
    \label{eq:ch6_track_parameters}
\end{align}
Here, $d_0$ is the signed transverse impact parameter, $z_0$ is the longitudinal coordinate at closest approach, $\phi_0$ and $\theta$ describe the direction of the momentum there, and $q/p$ is the charge-to-momentum ratio.
The fit also provides the covariance of these parameters.
The transverse impact-parameter significance, $d_0/\sigma(d_0)$, therefore expresses a displacement in units of its measurement uncertainty.
For longitudinal association to a vertex at $z_{\mathrm{PV}}$, the quantity $(z_0-z_{\mathrm{PV}})\sin\theta$ measures the longitudinal separation projected perpendicular to the track direction.
In the object selections below, this difference is denoted by $\Delta z_0\sin\theta$.

Prompt electrons and muons are expected to have small impact parameters, whereas tracks from $b$-hadron decays can have significant displacements.
Track requirements must therefore be chosen for their intended use.
The impact-parameter cuts used to select prompt leptons are not imposed on all tracks used for flavour tagging.
Likewise, tracks associated with a $\tau$ candidate undergo a dedicated classification to distinguish its charged decay products from nearby activity.

\section{Primary vertex}
\label{sec:ch6_primary_vertex}

Several proton--proton interactions can occur in the same bunch crossing, producing distinct collision vertices along the beam line.
Vertex reconstruction groups tracks compatible with a common production point and fits that point using the track parameters and their uncertainties~\cite{ATLAS2015Vertex}.
The vertex associated with the hard-scatter interaction provides the reference for prompt-object selection and for rejecting charged particles from pile-up.

In this analysis, a vertex must have at least two associated tracks with $\pT>500~\MeV$.
Among the reconstructed vertices, the primary vertex is chosen as the one with the largest sum of squared track transverse momenta~\cite{ATLAS2026LQTauNu},
\begin{align}
    v_{\mathrm{PV}}
    = \operatorname*{arg\,max}_{v}
      \sum_{i\in v} p_{\mathrm{T},i}^{\,2}.
    \label{eq:ch6_primary_vertex}
\end{align}
The quadratic weighting gives greater importance to energetic tracks, which are more characteristic of the hard scattering than of the additional soft interactions.
Events without a vertex satisfying these requirements are rejected.

The reconstructed primary vertex is an event-by-event position.
It differs from the \emph{beam spot}, which describes the spatial distribution of collision points over a period of data taking.
The beam-line reference and its uncertainty enter transverse impact-parameter measurements, while the reconstructed vertex identifies the collision to which a track is associated.
This association is used in lepton isolation, pile-up jet rejection and the track contribution to $\met$.
Hadronic $\tau$ reconstruction additionally uses a dedicated association of each candidate to its production vertex, as described in Section~\ref{sec:ch6_tau}.

\section{Electrons}
\label{sec:ch6_electrons}

An electron produces an inner-detector track and an electromagnetic shower in the calorimeter.
Bremsstrahlung in the material upstream of the calorimeter can alter the track curvature and distribute the energy over several nearby deposits.
Electron reconstruction therefore combines a track fit that accounts for bremsstrahlung with an electromagnetic \emph{supercluster}, which collects the main shower and compatible nearby clusters~\cite{ATLAS2019ElectronReco}.
Calorimeter clusters are built from connected cells with signals significant compared with the expected noise.
The spatial and energy measurements of the cluster are then matched to the extrapolated track.

Identification tests whether this combined signature is compatible with an electron.
The likelihood-based algorithm uses the lateral and longitudinal shower development, leakage into the hadronic calorimeter, track quality and track--cluster matching.
These quantities distinguish electrons from hadronic showers and other backgrounds.
The \texttt{TightLH} working point used in this analysis imposes a stringent requirement on the identification discriminant.
Its name refers to electron identification and is independent of the isolation requirement.

Isolation measures the additional activity around the candidate~\cite{ATLAS2019ElectronPerformance}.
An electron from a heavy-flavour decay is often accompanied by other particles in a jet, whereas an electron from a leptonic $W$ or $Z$ decay is more often isolated.
Track isolation sums the transverse momenta of nearby tracks after excluding the electron track.
Calorimeter isolation measures the surrounding transverse energy after subtracting the electron contribution and correcting for additional activity.
Variable-radius track isolation reduces the cone size as the electron momentum increases, limiting the loss of efficiency from nearby particles in energetic events.

Baseline electrons must satisfy $\pT>10~\GeV$ and $|\eta|<2.47$, excluding the calorimeter transition region $1.37<|\eta|<1.52$.
Candidates associated with problematic calorimeter regions are rejected.
They must pass \texttt{TightLH} identification and \texttt{Loose\_VarRad} isolation, together with
\begin{align}
    \frac{|d_0|}{\sigma(d_0)} < 5,
    \qquad
    |\Delta z_0\sin\theta| < 0.5~\mm.
    \label{eq:ch6_electron_ip}
\end{align}
These impact-parameter requirements suppress electrons from displaced decays and tracks inconsistent with the primary vertex.
Electrons selected for the control regions additionally satisfy \texttt{HighPtCaloOnly} isolation, which uses calorimeter activity to suppress backgrounds in the high-$\pT$ selection.
Their region-dependent momentum thresholds are specified with the event selections.

The electron energy is calibrated to correct for material upstream of the calorimeter, energy outside the reconstructed cluster and variations in detector response~\cite{ATLAS2019ElectronPerformance}.
Collision data, particularly $\zee$ events, constrain the remaining energy-scale and resolution differences between data and simulation.
Reconstruction, identification and isolation efficiencies are measured using tag-and-probe methods.
In this method, a well-identified electron selects a sample of resonance decays, and the second electron is used to measure the efficiency of the requirement under study after accounting for background.
The resulting corrections are applied for the working points used in the analysis.

\section{Muons}
\label{sec:ch6_muons}

Muons generally traverse the calorimeters with a relatively small energy deposit and leave measurements in the muon spectrometer.
Most selected muons are reconstructed by matching an inner-detector track to a muon-spectrometer track and fitting their measurements together, with the energy loss in the calorimeters included in the fit~\cite{ATLAS2021MuonID}.
The two tracking systems provide independent momentum information and a means of rejecting hadrons that decay in flight or produce misleading spectrometer signals.

Additional reconstruction categories recover muons for which a complete spectrometer track is unavailable.
A segment-tagged muon uses an inner-detector track matched to a short track segment in the spectrometer.
A calorimeter-tagged muon uses an inner-detector track with calorimeter deposits compatible with a minimum-ionising particle.
The latter category is relevant to the overlap-removal procedure because an electron and a calorimeter-tagged muon can be reconstructed from the same inner-detector track.

Muon identification combines requirements on detector hit content and on the consistency of the inner-detector and spectrometer measurements.
The \texttt{Loose} working point retains additional reconstruction categories to improve efficiency, while \texttt{Medium} imposes more restrictive quality requirements~\cite{ATLAS2021MuonID}.
Isolation further suppresses muons from heavy-flavour decays by requiring little nearby track and calorimeter activity.
As for electrons, the variable-radius track cone becomes smaller at higher $\pT$.

Baseline muons are required to have $\pT>10~\GeV$ and $|\eta|<2.5$, to satisfy \texttt{Loose} identification and \texttt{Loose\_VarRad} isolation, and to fulfil
\begin{align}
    \frac{|d_0|}{\sigma(d_0)} < 3,
    \qquad
    |\Delta z_0\sin\theta| < 0.5~\mm.
    \label{eq:ch6_muon_ip}
\end{align}
Muons in the control regions also pass \texttt{Medium} identification and \texttt{Tight\_VarRad} isolation.
The momentum calibration accounts for differences in detector alignment, magnetic-field description and material effects between data and simulation~\cite{ATLAS2023MuonCalibration}.
Resonance decays such as $\zmumu$ and $J/\psi\rightarrow\mu\mu$ constrain the momentum scale and resolution, while tag-and-probe measurements provide efficiency corrections for the reconstruction and selection requirements.

\section{Hadronic tau leptons}
\label{sec:ch6_tau}

\subsection{Decay signature and reconstruction}

The $\tau$-lepton has a proper lifetime of about $2.9\times10^{-13}$~s and is reconstructed through its decay products~\cite{ATLAS2022TauR22}.
Approximately 65\% of its decays are hadronic, producing a tau neutrino and a small number of hadrons~\cite{ATLAS2022TauR22}.
The dominant modes contain one or three charged hadrons, mainly pions, with possible additional neutral pions.
The charged hadrons produce tracks, while the photons from $\pi^0\rightarrow\gamma\gamma$ contribute electromagnetic energy deposits.
The boost of an energetic $\tau$-lepton collimates these particles, giving a narrow calorimeter signature with low charged-track multiplicity.

The reconstructed object represents the visible decay products and is denoted by $\tau_{\mathrm{had\text{-}vis}}$ when this distinction needs to be explicit.
The neutrino momentum is not included in its four-momentum.
For brevity, $\tau_{\mathrm{had}}$ denotes this reconstructed visible object in the following selections.
Consequently, the selected $\tau$ transverse momentum and a mass reconstructed from a $\tau$ candidate and a jet refer to visible quantities.
The tau-decay neutrino contributes to $\met$ together with the neutrino produced in the hard process.

Reconstruction starts from seed jets formed with the anti-$k_t$ algorithm with radius parameter $R=0.4$~\cite{Cacciari2008AntiKt,ATLAS2022TauR22}.
These seeds group topological calorimeter clusters, whose constituent cells are calibrated with the local hadronic calibration scheme~\cite{ATLAS2017TopoClusters}.
The scheme accounts for the different calorimeter responses to electromagnetic and hadronic energy deposits.
These calorimeter-based seeds are distinct from the particle-flow jets used for event categorisation in Section~\ref{sec:ch6_jets}.

A dedicated tau vertex association identifies the collision vertex most compatible with tracks within $\Delta R<0.2$ of the seed-jet direction.
For each vertex, the algorithm evaluates the fraction of the scalar track-$\pT$ sum associated with that vertex and chooses the largest fraction~\cite{ATLAS2022TauR22}.
This identifies the $\tau$ \emph{production} vertex and provides the reference for track association and impact parameters.
A secondary vertex reconstructed from three charged decay products instead estimates the displaced $\tau$ \emph{decay} position.
The two vertices therefore have different physical roles.

Tracks near the candidate are first required to satisfy geometrical and impact-parameter criteria.
A recurrent neural network (RNN) then classifies them as tau-decay tracks, conversion tracks, other nearby tracks or incorrectly reconstructed tracks~\cite{ATLAS2022TauR22}.
This classification is needed because an extra track from pile-up or a photon conversion can otherwise change the reconstructed decay multiplicity.
Candidates in this analysis must contain exactly one or three tracks classified as tau-decay tracks, with a total charge of $\pm1$ in units of the elementary charge.
They are referred to as one-prong and three-prong candidates, respectively.

\subsection{Identification and lepton rejection}

Quark- and gluon-initiated jets are produced much more frequently than hadronic $\tau$ decays and can occasionally have similar low-multiplicity signatures.
Tau identification exploits the narrower energy distribution, restricted track multiplicity and characteristic decay structure to suppress this background.
The RNN-based identifier combines sequences of associated tracks and calorimeter clusters with quantities describing the candidate as a whole~\cite{ATLAS2019TauRNN,ATLAS2022TauR22}.
The sequential inputs retain information about individual constituents and their correlations, including their momenta and angular separations.
Candidate-level inputs include the concentration of calorimeter energy around the axis, the mass of the track system and the amount of activity outside the core.
For three-prong decays, the significance of the transverse flight distance of the secondary vertex also contributes.
This jet-rejection network performs a different task from the track-classification network used during reconstruction.

Working points are defined by cuts on the identification score, with separate efficiency targets for one-prong and three-prong candidates.
Table~\ref{tab:ch6_tau_id} gives the targets used for the Run-2 reprocessing algorithms.
They describe identification conditional on reconstructing a suitable candidate and should not be interpreted as the total efficiency for a generated $\tau$-lepton to enter an analysis region.
The latter also depends on the decay mode, geometrical acceptance, reconstruction, lepton rejection and kinematic selection.
At high $\pT$, the charged decay products become more collimated and their tracks are harder to resolve individually.
This can reduce the reconstruction efficiency, particularly for three-prong candidates, even when the identification working point has a fixed target efficiency~\cite{ATLAS2022TauR22}.

\begin{table}[htbp]
    \centering
    \caption{Target efficiencies of the RNN tau-identification working points for reconstructed candidates matched to genuine hadronic $\tau$ decays, from Table~4 of Ref.~\cite{ATLAS2022TauR22}. The analysis uses \texttt{Loose} for baseline candidates and \texttt{Tight} for signal candidates.}
    \label{tab:ch6_tau_id}
    \begin{tabular}{lcc}
        \toprule
        Working point & One-prong & Three-prong \\
        \midrule
        \texttt{Loose}  & 85\% & 75\% \\
        \texttt{Medium} & 75\% & 60\% \\
        \texttt{Tight}  & 60\% & 45\% \\
        \bottomrule
    \end{tabular}
\end{table}

Electrons can also mimic one-prong $\tau$ decays because they produce a charged track and a concentrated electromagnetic shower.
A separate RNN electron-rejection discriminant uses tracking and calorimeter information, including shower development and the consistency of the measured energy with the track momentum~\cite{ATLAS2022TauR22}.
All baseline and signal candidates must satisfy its \texttt{Loose} working point.
Thus, tightening the tau jet-rejection requirement does not change the electron-rejection working point.
Candidates overlapping with identified electrons or muons are also removed according to Section~\ref{sec:ch6_overlap}.

\subsection{Energy calibration and analysis definition}

The tau energy calibration estimates the momentum of the visible decay products.
It combines calorimeter and tracking information, taking advantage of their different resolutions~\cite{ATLAS2022TauR22}.
The calorimeter-based estimate is combined with an estimate formed from reconstructed charged and neutral decay products.
A boosted regression tree then refines the response using information about the calorimeter shower, track multiplicity, reconstructed decay structure and pile-up.
The resulting multivariate tau energy scale is used for the candidate kinematics in this analysis.
It calibrates the visible object and does not recover the momentum carried by the tau-decay neutrino.

Baseline $\tau_{\mathrm{had}}$ candidates must have $\pT>20~\GeV$ and $|\eta|<2.5$, excluding $1.37<|\eta|<1.52$.
They must satisfy the one- or three-prong and unit-charge requirements, \texttt{Loose} RNN jet rejection and \texttt{Loose} RNN electron rejection.
Signal candidates additionally pass \texttt{Tight} RNN jet rejection.
The signal regions require $\pT^{\tau}>200~\GeV$, while some top-quark control and validation selections use a lower tau momentum threshold.
These thresholds belong to the region definitions and do not change the baseline collection used to count additional candidates.

Simulation is corrected for residual differences in the tau energy response and selection efficiencies using the corresponding ATLAS calibrations.
The analysis uses the electron-veto efficiency correction associated with \texttt{Medium} tau identification for candidates passing \texttt{Tight} tau identification, following the prescription recorded in the internal analysis note~\cite{LQInternalNote2026}.
% Internal note Section 5.4: no eVeto SF available specifically for Tight RNN;
% Medium-RNN eVeto SF used by agreement with the Tau CP group.
This prescription concerns the efficiency correction only.
The selected candidates still pass \texttt{Tight} jet rejection and \texttt{Loose} electron rejection.
Tau reconstruction and calibration are particularly relevant here because a change in the visible tau momentum affects both the tau selection and the reconstructed missing transverse momentum.

\section{Jets}
\label{sec:ch6_jets}

\subsection{Particle-flow reconstruction and clustering}

A high-energy quark or gluon produces a collimated collection of hadrons through parton showering and hadronisation.
A jet algorithm combines their measured signals into an object whose four-momentum approximates that of the resulting hadronic system.
The jets used for event selection are reconstructed from particle-flow constituents, combining inner-detector tracks and calorimeter clusters~\cite{ATLAS2017ParticleFlow}.
For charged particles, tracking can provide a more precise momentum measurement than calorimetry over part of the relevant momentum range.
The associated calorimeter contribution is subtracted when the track measurement is used, preventing the charged-particle momentum from being counted twice.
The remaining calorimeter signal contributes to the neutral component.
Vertex information also allows charged contributions from pile-up to be suppressed.

The particle-flow constituents are clustered with the anti-$k_t$ algorithm using $R=0.4$~\cite{Cacciari2008AntiKt}.
At each step, the algorithm compares pairwise distances and beam distances,
\begin{align}
    d_{ij}
    &= \min\!\left(p_{\mathrm{T},i}^{-2},p_{\mathrm{T},j}^{-2}\right)
       \frac{\Delta R_{y,ij}^{\,2}}{R^2},
    &
    d_{iB} &= p_{\mathrm{T},i}^{-2},
    \label{eq:ch6_antikt}\\
    \Delta R_{y,ij}
    &= \sqrt{(y_i-y_j)^2+(\phi_i-\phi_j)^2},
    &
    y &= \frac{1}{2}\ln\frac{E+p_z}{E-p_z}.
    \label{eq:ch6_rapidity}
\end{align}
Here, $y$ is rapidity, and the label $B$ denotes the beam distance rather than a physical beam constituent.
If a pairwise distance is smallest, the two four-momenta are combined.
If a beam distance is smallest, the corresponding object is declared a jet and removed from the clustering sequence.
The inverse-$\pT^2$ weighting causes energetic constituents to collect nearby soft radiation and gives isolated jets approximately circular boundaries of radius $R$.
The algorithm is infrared and collinear safe, meaning that adding an arbitrarily soft particle or replacing a particle by collinear fragments does not change the hard-jet configuration.

The distance $\Delta R_y$ uses rapidity, whereas $\Delta R$ in Eq.~\eqref{eq:delta_r} uses pseudorapidity.
They coincide for massless particles, but can differ for massive jets.
The distinction is retained explicitly below because the overlap-removal tools also use rapidity.

\subsection{Calibration, pile-up rejection and cleaning}

Jet calibration corrects the reconstructed four-momentum for additional energy from pile-up, the non-uniform detector response and energy losses~\cite{ATLAS2021JetCalibration}.
The first corrections remove the average pile-up contribution using the jet area and the transverse-momentum density of the event, followed by residual corrections for the pile-up dependence.
Simulation-based corrections then restore the average energy scale and direction.
Further corrections using jet properties reduce variations in response associated with the shower composition.
For collision data, residual differences from simulation are constrained with in situ momentum-balance measurements in events containing a well-measured photon, $Z$ boson or reference jet system.

The analysis retains jets with calibrated $\pT>20~\GeV$ and $|\eta|<2.5$.
Jets from additional collisions are suppressed with the neural-network jet-vertex tagger, NNJVT, at the \texttt{FixedEffPt} working point.
This tagger uses the association of charged tracks to vertices to distinguish jets from the hard scattering from pile-up jets.
The requirement is applied in the momentum and pseudorapidity range specified by the corresponding calibration recommendation.
% ANA_EXOT_2025_06_INT1, Table 5.4: NNJVT FixedEffPt; fJVT not applied.
The \texttt{LooseBad} jet-quality selection rejects detector signatures inconsistent with a jet from a collision, such as those arising from calorimeter noise or non-collision backgrounds.

The reconstructed jet multiplicity and momenta enter the event categorisation.
For resonant production, a jet can contain the quark from the leptoquark decay.
In non-resonant production, the accompanying jet probes a different part of the hard-scattering kinematics.
The requirements that distinguish these topologies are introduced in Chapter~\ref{chapter:event_selection}.
The inclusive jet collection defined here is also the starting point for identifying $b$-tagged jets.

\section{Jet flavour tagging}
\label{sec:ch6_btag}

Flavour tagging distinguishes jets containing heavy-flavour hadrons from jets initiated by light quarks or gluons.
The characteristic signature of a $b$-jet arises from the lifetime of the $b$-hadron.
Its mean flight distance is
\begin{align}
    \langle L\rangle = \beta\gamma c\tau_b,
    \label{eq:ch6_b_flight}
\end{align}
where $\tau_b$ is its proper lifetime and $\beta\gamma$ is its boost.
With a typical lifetime of about $1.5$~ps, a boosted $b$-hadron can travel several millimetres before decaying~\cite{ATLAS2026GN2}.
Its charged decay products therefore have displaced impact parameters and may form a secondary vertex.
The subsequent decay of a charm hadron can produce a further displaced vertex.
These spatial correlations, together with the decay multiplicity and track momenta, distinguish $b$-jets from most light-flavour jets.
Charm jets also contain displaced decays and are consequently a separate background to $b$-tagging.

This analysis uses the GN2v01 flavour-tagging algorithm.
GN2 processes the kinematics and tracking information of the jet constituents with a transformer-based neural network~\cite{ATLAS2026GN2}.
The input includes the jet momentum and direction, track momenta, impact parameters and information about detector hits.
The transformer allows each track representation to depend on the other tracks in the same jet, making correlated displacements and common decay origins available to the classifier.
During training, auxiliary tasks predict track origins and whether pairs of tracks share a production vertex.
These tasks encourage the network to learn the decay structure that carries the flavour information.
GN2 therefore does not require a list of explicitly reconstructed secondary vertices as an input to a separate final discriminant.

The network produces outputs for the $b$-, $c$-, hadronic-tau and light-jet hypotheses.
They are combined into a $b$-tagging discriminant of the form
\begin{align}
    D_b = \ln\!\left[
        \frac{p_b}
        {f_c p_c + f_\tau p_\tau +(1-f_c-f_\tau)p_u}
    \right],
    \label{eq:ch6_gn2_discriminant}
\end{align}
where $p_b$, $p_c$, $p_\tau$ and $p_u$ are the respective classifier outputs~\cite{ATLAS2026GN2}.
The fixed coefficients $f_c$ and $f_\tau$ tune the balance between rejecting different background flavours.
They are properties of the discriminant and do not represent measured flavour fractions in the selected data.
A jet is tagged when $D_b$ exceeds the threshold defining the chosen working point.

The nominal selection uses the fixed-cut working point corresponding to an inclusive $b$-jet efficiency of 85\% in simulated $\ttbar$ events.
This efficiency defines the reference operating point and is not a constant efficiency for every jet momentum or every signal model.
For a background flavour $f$, the rejection is $1/\epsilon_f$, where $\epsilon_f$ is the probability of that flavour passing the same requirement.
Both the tagging and mistagging efficiencies depend on the jet kinematics and are corrected using the appropriate flavour-dependent calibrations.

These corrections affect events with tagged and untagged jets.
For example, consider a single jet with simulated tagging efficiency $\epsilon_f^{\mathrm{MC}}$ and scale factor $\mathrm{SF}_f=\epsilon_f^{\mathrm{data}}/\epsilon_f^{\mathrm{MC}}$.
The correction for the tagged outcome is $\mathrm{SF}_f$, while that for the untagged outcome follows from the complementary probability,
\begin{align}
    w_f^{\mathrm{untagged}}
    = \frac{1-\epsilon_f^{\mathrm{data}}}
           {1-\epsilon_f^{\mathrm{MC}}}
    = \frac{1-\mathrm{SF}_f\epsilon_f^{\mathrm{MC}}}
           {1-\epsilon_f^{\mathrm{MC}}}.
    \label{eq:ch6_btag_untagged}
\end{align}
This relation explains why a $b$-tagging calibration changes the predicted population of the zero-tag category as well as that of the tagged category.
The analysis applies the corresponding calibration through the flavour-tagging tools.

The number of tagged jets separates the zero-$b$-tag and one-$b$-tag signal categories.
The former can receive signal events when a $b$-jet is outside acceptance or fails the tag.
It also has a direct signal contribution from $U_1\rightarrow s\tau$ when $\beta_L^{23}$ is sizeable.
The zero-tag category therefore probes part of the signal flavour structure in its own right.
It cannot be interpreted solely as a correction for imperfect $b$-tagging.
Events with additional $b$-tagged jets are useful for constraining top-quark backgrounds, as discussed in Chapter~\ref{chapter:background_estimation}.

\section{Overlap removal}
\label{sec:ch6_overlap}

Object reconstruction algorithms operate on partially shared detector information.
An electron shower can also be reconstructed as a jet, and a hadronic $\tau$ candidate is itself seeded by a calorimeter jet.
Counting all such candidates independently would distort object multiplicities and assign more than one interpretation to the same activity.
An ordered overlap-removal procedure selects the objects used for event selection.
Only objects surviving an earlier step are considered in subsequent steps.

Table~\ref{tab:ch6_overlap} gives the sequence reported in the analysis paper~\cite{ATLAS2026LQTauNu}.
The matching distances use $\Delta R_y$, defined in Eq.~\eqref{eq:ch6_rapidity}, in accordance with the rapidity convention of the overlap-removal tools.
An inner electron--jet overlap is resolved in favour of the electron because the jet may be the electron shower reconstructed a second time.
If a jet survives this step, an electron within the larger outer cone is removed because it is less consistent with an isolated electron.
For muons, the jet track multiplicity helps distinguish calorimeter activity associated with the muon from a genuine hadronic jet containing a muon.
The final tau--jet step retains the tau interpretation of an overlapping jet.

\begin{table}[htbp]
    \centering
    \small
    \caption{Ordered overlap-removal sequence as described in Ref.~\cite{ATLAS2026LQTauNu}. Each row is applied to candidates surviving the preceding rows. The muon--jet matching description and an additional electron--electron step in the internal note require the version checks recorded in the accompanying draft notes.}
    \label{tab:ch6_overlap}
    \renewcommand{\arraystretch}{1.16}
    \begin{tabular}{clp{8.1cm}}
        \toprule
        Step & Rejected object & Condition \\
        \midrule
        1 & $\tau_{\mathrm{had}}$ & Within $\Delta R_y<0.2$ of an electron. \\
        2 & $\tau_{\mathrm{had}}$ & Within $\Delta R_y<0.2$ of a muon. \\
        3 & Muon & Calorimeter-tagged and sharing an inner-detector track with an electron. \\
        4 & Electron & Sharing an inner-detector track with a muon. \\
        5 & Jet & The closest jet within $\Delta R_y<0.2$ of an electron. \\
        6 & Electron & Within $\Delta R_y<0.4$ of a surviving jet. \\
        7 & Jet & Within $\Delta R_y<0.4$ of a muon and having at most two associated tracks, as specified in the paper. \\
        8 & Muon & Within $\Delta R_y<0.4$ of a surviving jet. \\
        9 & Jet & Within $\Delta R_y<0.2$ of a surviving $\tau_{\mathrm{had}}$ candidate. \\
        \bottomrule
    \end{tabular}
\end{table}

Overlap removal for object counting must also be distinguished from the treatment of shared detector signals in the missing-momentum calculation.
The latter uses associations between tracks, calorimeter signals and reconstructed objects to account for their momenta consistently.
It is described in the following section.

\section{Missing transverse momentum}
\label{sec:ch6_met}

The transverse momenta of the incoming partons are small compared with the momentum scale of the hard scattering.
Momentum conservation therefore predicts an approximately balanced final state in the transverse plane when all particles are measured.
Neutrinos escape detection and produce an imbalance in the measured visible momentum.
In the signal considered here, both the hard-process neutrino and the tau-decay neutrino can contribute to this imbalance.
Because they contribute as vectors, the measured imbalance depends on their directions as well as on their individual momenta.

The reconstructed missing transverse momentum is the negative vector sum of the calibrated visible contributions~\cite{ATLAS2025METPerformance},
\begin{align}
    \vec p_{\mathrm{T}}^{\,\mathrm{miss}}
    &= -\left(
        \sum_{e}\vec p_{\mathrm{T}}^{\,e}
        + \sum_{\mu}\vec p_{\mathrm{T}}^{\,\mu}
        + \sum_{\tau_{\mathrm{had}}}\vec p_{\mathrm{T}}^{\,\tau}
        + \sum_{\mathrm{jets}}\vec p_{\mathrm{T}}^{\,j}
        + \vec p_{\mathrm{T}}^{\,\mathrm{soft}}
       \right),
    \label{eq:ch6_met_vector}\\
    \met &= \left|\vec p_{\mathrm{T}}^{\,\mathrm{miss}}\right|.
    \label{eq:ch6_met_magnitude}
\end{align}
The object sums form the \emph{hard term} and contain the contributions accepted by the missing-momentum reconstruction.
The \emph{soft term}, $\vec p_{\mathrm{T}}^{\,\mathrm{soft}}$, accounts for additional charged activity associated with the primary vertex.
No separate photon collection is supplied in this analysis, so Eq.~\eqref{eq:ch6_met_vector} contains the electron, muon, tau and jet collections used by its reconstruction configuration.

The baseline lepton and tau selections provide the candidate inputs.
An association map records which detector signals belong to each object, and the reconstruction resolves shared contributions before forming the sum~\cite{ATLAS2025METPerformance}.
It can therefore treat partial overlaps and muon energy deposits within jets in more detail than the geometrical object-removal procedure.
The jet term is subject to its own missing-momentum working point and should not be identified with the list of central jets counted in the event selection.
In particular, the calculation is not a sum of only the high-$\pT$ objects satisfying the final signal-region requirements.

The analysis uses the track soft term and the \texttt{Tight} missing-momentum operating point.
The soft term includes suitable primary-vertex tracks not assigned to retained hard-object contributions, with additional quality and overlap requirements.
Restricting it to primary-vertex tracks suppresses pile-up contamination because tracks can be associated with their production vertices.
Soft neutral particles without an accepted hard-object contribution are not measured by this term.
This is a limitation of its response, balanced by its reduced sensitivity to pile-up~\cite{ATLAS2025METPerformance}.
The \texttt{Tight} operating point controls the jet contribution to the calculation and is independent of the \texttt{Tight} tau-identification requirement.

Large $\met$ can also arise when visible energy is mismeasured, lost in uninstrumented material or affected by detector noise.
For example, underestimating a jet momentum can produce an apparent missing-momentum vector along the jet direction.
Jet-quality requirements and selections using the azimuthal separation between jets and $\vec p_{\mathrm{T}}^{\,\mathrm{miss}}$ suppress these contributions.
Their precise event-level definitions are introduced in Chapter~\ref{chapter:event_selection}.
Energy-scale and resolution variations of the hard objects are propagated to $\vec p_{\mathrm{T}}^{\,\mathrm{miss}}$, while separate uncertainties describe the track soft term.
This treatment preserves the relationship between the visible-object kinematics and the momentum imbalance used to select the signal.
